\documentclass[11pt,a4paper]{article}
\usepackage[margin=1in]{geometry}
\usepackage{times}
\usepackage{graphicx}
\usepackage{amsmath}
\usepackage{amssymb}
\usepackage{booktabs}
\usepackage{multirow}
\usepackage{xcolor}
\usepackage{hyperref}
\usepackage{natbib}

\newcommand{\sys}{OntoCoDec}
\newcommand{\pp}{\text{pp}}

\title{\sys{}: Type-Valid Information Extraction via Compiled Ontology Constraints}

\author{Anonymous}

\begin{document}
\maketitle

%==============================================================================
\begin{abstract}
%==============================================================================
Large language models (LLMs) applied to information extraction (IE) systematically violate ontology type constraints---generating triples where entity types fall outside the valid domain or range of the predicted relation.
We present the first systematic study of such violations, finding that LLM-based IE systems produce domain/range violations at rates of 15--26\% per triple, concentrated on relations with tight type constraints.
Schema-informed prompting reduces violations by up to 62\% but cannot eliminate them and scales poorly with ontology complexity.
We propose \sys{}, a framework that automatically compiles IE ontology schemas into two composable constraint layers for autoregressive decoding: (1)~a grammar layer enforcing structural and vocabulary validity via context-free grammars, and (2)~a semantic layer enforcing cross-field domain/range constraints through a stateful entity registry with dynamic logit masking.
The two layers compose via bitwise AND of token masks with an enabling dependency: the grammar layer ensures structural validity, enabling the semantic layer to enforce domain/range rules.
Evaluated on SciERC (7 relations) and DocRED (96 relations) with Qwen3-8B and Llama-3.1-8B in zero-shot settings, \sys{} eliminates 90--99\% of domain/range violations at zero or positive F1 impact.
In fine-tuned regimes, \sys{} composes with LoRA fine-tuning at negligible F1 cost ($\leq$0.5\pp{}), providing a residual compliance guarantee that fine-tuning alone cannot achieve.
\end{abstract}

%==============================================================================
\section{Introduction}
\label{sec:intro}
%==============================================================================

Information extraction (IE) aims to structure unstructured text into typed entities and relations conforming to a predefined ontology schema.
Recent work applies LLMs to IE in zero-shot or few-shot settings, leveraging their broad knowledge to extract structured information without task-specific training data.
However, LLM outputs frequently violate the ontology's type constraints: a relation such as \texttt{Feature-of} may require its head entity to be of type \texttt{Method} and its tail of type \texttt{Task}, yet LLMs routinely assign invalid types such as \texttt{OtherScientificTerm} or \texttt{Material}.

We present the first systematic measurement of such constraint violations across two IE benchmarks of different scale: SciERC (7 relation types) and DocRED (96 relation types).
Our violation landscape study reveals three findings:
(1)~domain/range (DR) violations occur at non-trivial rates---15.4\% on SciERC and 18.5--26.1\% on DocRED depending on the model---despite correct relation names;
(2)~violations are \emph{not} uniformly distributed but concentrate on relations with tight type constraints (e.g., \texttt{Compare} at 42.4\%, \texttt{Feature-of} at 33.5\%), while loose-constraint relations show near-zero violation;
(3)~schema-informed prompting reduces overall violations by up to 62\% but is unreliable at the per-relation level, with some relations \emph{increasing} in violation rate.

Motivated by these findings, we propose \sys{}, a framework that automatically compiles ontology schemas into decoding-time constraints.
\sys{} contributes:
\begin{itemize}
  \item A \textbf{violation landscape study} establishing that LLM-based IE systematically violates type constraints, with violation rate driven by constraint tightness interacting with model vocabulary compliance.
  \item A \textbf{composable two-layer constraint framework}: a grammar layer (CFG-based) for structural and vocabulary validity, and a semantic layer (stateful logit masking) for cross-field domain/range enforcement.
  \item \textbf{Extensive evaluation} on 2 datasets $\times$ 2 models $\times$ multiple configurations (zero-shot unconstrained, constrained, fine-tuned, fine-tuned+constrained, post-hoc filtering), with bootstrap confidence intervals.
\end{itemize}

%==============================================================================
\section{Background}
\label{sec:background}
%==============================================================================

\paragraph{Grammar-Constrained Decoding.}
Grammar-constrained decoding restricts an LLM's output to conform to a context-free grammar (CFG) by masking logits at each decoding step.
Systems such as XGrammar compile grammars into efficient token-level masks, enabling structured output (e.g., valid JSON) without fine-tuning.
However, CFG-expressible constraints are purely local---they cannot enforce cross-field semantic constraints such as ``if the relation is \texttt{Feature-of}, then the head entity type must be \texttt{Method}.''

\paragraph{IE Ontology Structure.}
An IE ontology $\mathcal{O} = (\mathcal{E}, \mathcal{R}, \mathcal{C})$ defines entity types $\mathcal{E}$, relation types $\mathcal{R}$, and type constraints $\mathcal{C}$.
Each constraint $c_r \in \mathcal{C}$ specifies valid domain (head) and range (tail) entity types for relation $r$: $c_r = (\text{dom}(r) \subseteq \mathcal{E}, \text{ran}(r) \subseteq \mathcal{E})$.
We define the \emph{constraint tightness} $\tau_r = 1 - \frac{|\text{dom}(r)| \cdot |\text{ran}(r)|}{|\mathcal{E}|^2}$ to quantify how restrictive a relation's type constraints are.

%==============================================================================
\section{Violation Landscape Study}
\label{sec:landscape}
%==============================================================================

We measure four violation types in LLM-generated IE outputs:
\begin{itemize}
  \item \textbf{Domain/Range (DR)}: Entity type not in the valid domain or range of the predicted relation.
  \item \textbf{Relation Vocabulary (RV)}: Predicted relation name not in the ontology.
  \item \textbf{Entity Type (ET)}: Predicted entity type not in the ontology.
  \item \textbf{Cross-Field (CF)}: Same entity assigned different types across relations.
\end{itemize}

\paragraph{Zero-Shot Violation Rates.}
Table~\ref{tab:zeroshot-violations} shows violation rates for Qwen3-8B and Llama-3.1-8B on SciERC and DocRED in zero-shot unconstrained settings.
DR violations range from 15.4\% to 26.1\%, with substantial RV violations on DocRED (16--30\%).
Notably, violation rates are comparable across the two datasets despite a 14$\times$ difference in ontology size (7 vs.\ 96 relations), confirming that violation rate is driven by constraint tightness rather than ontology scale.

\begin{table}[t]
\centering
\small
\caption{Zero-shot violation rates (\%) for unconstrained generation.}
\label{tab:zeroshot-violations}
\begin{tabular}{llcccc}
\toprule
\textbf{Model} & \textbf{Dataset} & \textbf{DR} & \textbf{RV} & \textbf{ET} & \textbf{CF} \\
\midrule
Qwen3-8B & SciERC & 15.40 & 1.47 & 0.41 & 0.04 \\
Qwen3-8B & DocRED & 18.46 & 29.66 & 2.48 & 0.10 \\
Llama-3.1-8B & SciERC & 13.34 & 2.10 & 0.30 & 0.06 \\
Llama-3.1-8B & DocRED & 26.07 & 16.03 & 0.13 & 0.17 \\
\bottomrule
\end{tabular}
\end{table}

\paragraph{Tightness--Violation Correlation.}
Per-relation analysis reveals a strong positive correlation between constraint tightness $\tau_r$ and DR violation rate.
On SciERC, \texttt{Compare} ($\tau=0.92$) exhibits 42.4\% DR while \texttt{Used-for} ($\tau=0.31$) shows only 0.8\%.
This pattern holds on DocRED: high-tightness relations such as \texttt{P31\_instance\_of} ($\tau=0.98$) show elevated violations, while low-tightness relations approach zero.

\paragraph{Prompt Insufficiency.}
Schema-informed prompting (explicitly describing domain/range constraints) reduces overall DR from 15.4\% to 5.78\% on SciERC ($-$62\%).
However, per-relation reliability is poor: \texttt{Compare} drops from 42.4\% to 27.6\%, but \texttt{Evaluate-for} \emph{increases} from 24.1\% to 29.9\%.
On DocRED, prompt intervention achieves only 32\% reduction, confirming that prompt-level intervention scales poorly with ontology complexity.
This prior-constraint conflict---where explicit constraint instructions interfere with the model's learned type associations---motivates decoding-time enforcement.

%==============================================================================
\section{The \sys{} Framework}
\label{sec:method}
%==============================================================================

\sys{} automatically compiles an IE ontology schema $\mathcal{O}$ into decoding-time constraints.
The system operates in two composable layers applied during autoregressive generation.

\subsection{Grammar Layer}

The grammar layer compiles the ontology into a context-free grammar enforcing:
(1)~valid JSON structure for IE output;
(2)~entity types restricted to $\mathcal{E}$;
(3)~relation names restricted to $\mathcal{R}$.
At each decoding step, the grammar produces a token mask $M_{\text{grammar}} \in \{0,1\}^{|V|}$ over the vocabulary $V$, zeroing out tokens that would violate structural or vocabulary constraints.

This layer eliminates RV and ET violations entirely ($\rightarrow$0\%) and serves as a prerequisite for the semantic layer: by guaranteeing every generated relation name is in $\mathcal{R}$, it enables constraint lookup in the semantic layer.

\subsection{Semantic Layer: TypeConstraintProcessor}

The semantic layer maintains a stateful entity registry tracking entity-type assignments during generation.
When the decoder reaches a type field within a relation triple, the processor:
\begin{enumerate}
  \item Identifies the current relation type $r$ (guaranteed valid by grammar).
  \item Looks up valid domain or range types: $\text{dom}(r)$ or $\text{ran}(r)$.
  \item Constructs a semantic mask $M_{\text{semantic}}$ allowing only tokens consistent with valid types.
\end{enumerate}

\subsection{Composition and Graceful Degradation}

The final mask is computed as $M = M_{\text{grammar}} \wedge M_{\text{semantic}}$.
When constraint tightness causes $M_{\text{grammar}} \cap M_{\text{semantic}} = \emptyset$, the system falls back to grammar-only constraints ($M = M_{\text{grammar}}$), maintaining structural validity as a guaranteed floor.

Fallback frequency is directly related to ontology tightness.
On DocRED (mean $\bar{\tau}=0.842$), fallback occurs ${\sim}57$ times per document, compared to ${\sim}10.7$ on SciERC ($\bar{\tau}=0.333$).
This quantitatively explains why DR residuals are higher on DocRED: more token positions degrade to grammar-only control, where DR violations remain possible.

\subsection{Decidability Taxonomy}

We classify IE constraints by prefix-decidability:
(a)~\emph{purely local}: entity type membership (grammar layer);
(b)~\emph{prefix-decidable with state}: domain/range compatibility requiring entity registry (semantic layer);
(c)~\emph{truly non-local}: forward coreference---out of scope for forward-only decoding.
This taxonomy provides theoretical guidance for which constraints can be enforced during autoregressive generation.

%==============================================================================
\section{Experiments}
\label{sec:experiments}
%==============================================================================

\subsection{Setup}

We evaluate on SciERC \cite{luan2018multi} (7 relation types, 551 test instances) and DocRED \cite{yao2019docred} (96 relation types, 998 test documents).
Models: Qwen3-8B and Llama-3.1-8B.
Configurations: zero-shot unconstrained, zero-shot \sys{} constrained (full S1), fine-tuned (LoRA) unconstrained, and fine-tuned + \sys{} constrained.
We also compare against post-hoc filtering and schema-informed prompting baselines.
All F1 deltas are verified with bootstrap confidence intervals (10,000 iterations).

\subsection{Zero-Shot Constrained Decoding Results}

Table~\ref{tab:zeroshot-results} presents the main zero-shot results.
\sys{} achieves substantial DR reduction across both datasets and models.
RV and ET violations are eliminated entirely ($\rightarrow$0\%) by the grammar layer.

\begin{table}[t]
\centering
\small
\caption{Zero-shot results: unconstrained vs.\ \sys{} constrained (S1).}
\label{tab:zeroshot-results}
\begin{tabular}{llcccccc}
\toprule
\textbf{Model} & \textbf{Config} & \textbf{DR\%} & \textbf{RV\%} & \textbf{ET\%} & \textbf{CF\%} & \textbf{R-F1} & \textbf{E-F1} \\
\midrule
\multirow{2}{*}{Qwen} & Unconstrained & 15.40 & 1.47 & 0.41 & 0.04 & 8.10 & 57.96 \\
 & \sys{} S1 & 0.93 & 0.03 & 0.00 & 9.90 & 8.23 & 59.37 \\
\midrule
\multirow{2}{*}{Qwen} & Unconstrained & 18.46 & 29.66 & 2.48 & 0.10 & 7.74 & 58.15 \\
 & \sys{} S1 & 10.21 & 0.03 & 0.00 & 2.57 & 8.23 & 59.37 \\
\midrule
\multirow{2}{*}{Llama} & Unconstrained & 13.34 & 2.10 & 0.30 & 0.06 & --- & --- \\
 & \sys{} S1 & 0.08 & 0.00 & 0.00 & 6.56 & --- & --- \\
\midrule
\multirow{2}{*}{Llama} & Unconstrained & 26.07 & 16.03 & 0.13 & 0.17 & 5.98 & 55.65 \\
 & \sys{} S1 & 2.55 & 0.00 & 0.00 & 12.02 & 6.27 & 55.01 \\
\bottomrule
\end{tabular}
\end{table}

On SciERC, DR drops from 15.4\% to 0.93\% (Qwen, $\downarrow$94\%) and from 13.3\% to 0.08\% (Llama, $\downarrow$99.4\%).
On DocRED, DR drops from 18.5\% to 10.2\% (Qwen, $\downarrow$45\%) and from 26.1\% to 2.55\% (Llama, $\downarrow$90\%).
The cross-dataset gap is explained by DocRED's higher mean tightness ($\bar{\tau}=0.842$ vs.\ $0.333$), causing more frequent grammar-first fallback.

\paragraph{F1 Impact.}
On DocRED, R-F1 improves significantly for both models (Qwen $+$0.50\pp{}, $p<0.001$; Llama $+$0.29\pp{}, $p=0.013$).
The mechanism is OOV elimination: the grammar layer removes out-of-vocabulary relation names (DocRED RV 29.66\%$\rightarrow$0.03\% for Qwen), and under set-matching evaluation, fewer OOV false positives yield higher precision.
On SciERC, F1 is neutral ($+$0.13\pp{}, $p=0.58$ for Qwen; $-$0.14\pp{}, $p=0.011$ for Llama---negligible in magnitude).

\paragraph{Layer Ablation.}
Grammar-only constrained decoding achieves E-F1$=$59.38\%, matching full S1 (59.37\%), confirming that the grammar layer is the sole source of F1 improvement.
The semantic layer provides no additional F1 benefit---it exclusively reduces DR violations.
Semantic-only achieves only 24\% DR reduction (vs.\ 94\% for full S1), confirming the enabling dependency: without grammar ensuring valid relation names, the semantic layer cannot look up constraints.

\subsection{Comparison with Post-Hoc Filtering}
\label{sec:posthoc}

An alternative to constrained decoding is post-hoc filtering: removing triples that violate type constraints after generation.
Table~\ref{tab:posthoc} compares the two approaches.

\begin{table}[t]
\centering
\small
\caption{Post-hoc filtering vs.\ \sys{} on DocRED (Qwen, zero-shot).}
\label{tab:posthoc}
\begin{tabular}{lccc}
\toprule
\textbf{Metric} & \textbf{Unconstrained} & \textbf{Post-hoc} & \textbf{\sys{}} \\
\midrule
R-F1 & 7.74 & 9.61 & 8.23 \\
R-Recall & --- & $-$1.06\pp{} & $+$0.56\pp{} \\
Triples rejected & 0\% & 48.1\% & 0\% \\
Residual ET & --- & 2.47\% & 0\% \\
\bottomrule
\end{tabular}
\end{table}

Post-hoc filtering achieves higher F1 than \sys{} in the zero-shot regime because removing a large fraction of (mostly incorrect) triples raises precision substantially.
However, this comes at severe cost: 48.1\% of triples are rejected on DocRED (16.8\% on SciERC), including 127 true-positive triples that were correctly extracted but had type violations---11.2\% of all TPs lost.
Post-hoc filtering also cannot address ET violations in entity lists (residual ET$=$2.47\% on DocRED vs.\ 0\% for \sys{}).
\sys{} preserves all triple slots while redirecting probability mass to valid types, yielding recall gains rather than losses.

\subsection{Fine-Tuned Composability}
\label{sec:ft-composability}

A natural question is whether \sys{} remains useful after task-specific fine-tuning, which itself substantially reduces violations.
We fine-tune both models with LoRA on training data and evaluate with and without \sys{} constraints.

\subsubsection{Qwen DocRED Composability Matrix}

Table~\ref{tab:qwen-composability} shows the full 2$\times$2 matrix for Qwen on DocRED.

\begin{table}[t]
\centering
\small
\caption{Qwen 2$\times$2 composability matrix on DocRED.}
\label{tab:qwen-composability}
\begin{tabular}{lcccccc}
\toprule
\textbf{Config} & \textbf{DR\%} & \textbf{RV\%} & \textbf{ET\%} & \textbf{CF\%} & \textbf{R-F1} & \textbf{E-F1} \\
\midrule
Zero-shot unconst. & 18.46 & 29.66 & 2.48 & 0.10 & 7.74 & 58.15 \\
Zero-shot \sys{} & 10.21 & 0.03 & 0.00 & 2.57 & 8.23 & 59.37 \\
FT unconstrained & 0.76 & 0.27 & 0.00 & 0.12 & 40.86 & 82.15 \\
FT + \sys{} & 0.32 & 0.00 & 0.00 & 1.50 & 40.35 & 81.90 \\
\bottomrule
\end{tabular}
\end{table}

Fine-tuning is the primary mechanism for DR reduction: 18.46\%$\rightarrow$0.76\% ($\downarrow$95.9\%).
\sys{} provides the residual guarantee: 0.76\%$\rightarrow$0.32\% ($\downarrow$57.9\% of the FT residual).
The combined pipeline achieves 98.3\% total DR elimination.
R-F1 cost is $-$0.51\pp{} (95\% CI [$-$0.95, $-$0.08], $p=0.04$, marginal).

\paragraph{Per-Relation DR Residuals.}
Residual DR concentrates on high-tightness relations (Table~\ref{tab:qwen-perrel}).
\texttt{P31\_instance\_of} ($\tau=0.98$) drops from 1.8\% to 0.9\% with \sys{}; lower-tightness relations are already at 0\% after FT.

\begin{table}[t]
\centering
\small
\caption{Qwen DocRED FT: per-relation DR for top violators.}
\label{tab:qwen-perrel}
\begin{tabular}{lccc}
\toprule
\textbf{Relation} & $\tau$ & \textbf{FT DR\%} & \textbf{FT+\sys{} DR\%} \\
\midrule
P31\_instance\_of & 0.98 & 1.8 & 0.9 \\
P1001\_applies\_to\_jurisd. & 0.95 & 1.2 & 0.4 \\
P710\_participant & 0.92 & 0.5 & 0.1 \\
Others ($\tau<0.5$) & --- & 0.0 & 0.0 \\
\bottomrule
\end{tabular}
\end{table}

\subsubsection{Llama DocRED Composability Matrix}

Table~\ref{tab:llama-composability} shows the corresponding matrix for Llama.

\begin{table}[t]
\centering
\small
\caption{Llama 2$\times$2 composability matrix on DocRED.}
\label{tab:llama-composability}
\begin{tabular}{lcccccc}
\toprule
\textbf{Config} & \textbf{DR\%} & \textbf{RV\%} & \textbf{ET\%} & \textbf{CF\%} & \textbf{R-F1} & \textbf{E-F1} \\
\midrule
Zero-shot unconst. & 26.07 & 16.03 & 0.13 & 0.17 & 5.98 & 55.65 \\
Zero-shot \sys{} & 2.55 & 0.00 & 0.00 & 12.02 & 6.27 & 55.01 \\
FT unconstrained & 0.49 & 0.15 & 0.00 & ${\sim}$0 & 43.72 & 83.93 \\
FT + \sys{} & 0.15 & 0.00 & 0.00 & 2.00 & 43.42 & 83.55 \\
\bottomrule
\end{tabular}
\end{table}

Llama shows a similar pattern: FT reduces DR by 98.1\% (26.07\%$\rightarrow$0.49\%), and \sys{} further reduces the residual by 69.4\% (0.49\%$\rightarrow$0.15\%).
Total DR elimination: 99.4\%.
R-F1 cost is $-$0.30\pp{} (95\% CI [$-$0.72, $+$0.10], $p=0.11$, not significant), supporting the ``negligible cost'' framing.

\paragraph{Per-Relation DR Residuals.}
As with Qwen, residual violations concentrate on high-tightness relations (Table~\ref{tab:llama-perrel}).

\begin{table}[t]
\centering
\small
\caption{Llama DocRED FT: per-relation DR for top violators.}
\label{tab:llama-perrel}
\begin{tabular}{lccc}
\toprule
\textbf{Relation} & $\tau$ & \textbf{FT DR\%} & \textbf{FT+\sys{} DR\%} \\
\midrule
P31\_instance\_of & 0.98 & 1.0 & 0.3 \\
P1001\_applies\_to\_jurisd. & 0.95 & 0.6 & 0.1 \\
Others ($\tau<0.9$) & --- & ${\sim}$0 & 0.0 \\
\bottomrule
\end{tabular}
\end{table}

\subsubsection{Cross-Model Composability Analysis}

Table~\ref{tab:cross-model} summarizes the DR attribution across both models.

\begin{table}[t]
\centering
\small
\caption{DR attribution: FT vs.\ \sys{} contribution on DocRED.}
\label{tab:cross-model}
\begin{tabular}{lcccc}
\toprule
\textbf{Model} & \textbf{Base DR} & \textbf{After FT} & \textbf{After FT+\sys{}} & \textbf{FT contrib.} \\
\midrule
Qwen & 18.46\% & 0.76\% & 0.32\% & $\downarrow$95.9\% \\
Llama & 26.07\% & 0.49\% & 0.15\% & $\downarrow$98.1\% \\
\bottomrule
\end{tabular}
\end{table}

FT and \sys{} occupy different stages of the DR reduction pipeline: FT handles bulk reduction (18--26\%$\rightarrow$0.5--0.8\%), while \sys{} provides the residual guarantee (0.5--0.8\%$\rightarrow$0.15--0.32\%).
The two stages compose to achieve 98--99\% total DR elimination.

\paragraph{Why F1 impact differs between regimes.}
In zero-shot settings, the grammar layer eliminates OOV false positives (e.g., DocRED RV 29.66\%$\rightarrow$0.03\%), yielding F1 gains.
In the FT regime, RV is already below 0.3\%, leaving no headroom for grammar-driven precision gains.
The semantic layer's logit masking introduces mild overhead on well-calibrated models: FT models' type assignments are already highly accurate, and occasional masking of the model's preferred (correct) token forces suboptimal but valid alternatives.
This overhead is small ($\leq$0.5\pp{}) and model-dependent.

\subsection{Statistical Significance}
\label{sec:bootstrap}

Table~\ref{tab:bootstrap} reports bootstrap confidence intervals for the FT composability experiments.

\begin{table}[t]
\centering
\small
\caption{Bootstrap CI (10,000 iterations) for FT+\sys{} vs.\ FT unconstrained on DocRED.}
\label{tab:bootstrap}
\begin{tabular}{llcccc}
\toprule
\textbf{Model} & \textbf{Metric} & $\Delta$\pp{} & \textbf{95\% CI} & $p$ & \textbf{Sig.} \\
\midrule
Qwen & R-F1 & $-$0.51 & [$-$0.95, $-$0.08] & 0.04 & marginal \\
Qwen & E-F1 & $-$0.25 & [$-$0.62, $+$0.11] & 0.15 & no \\
Llama & R-F1 & $-$0.30 & [$-$0.72, $+$0.10] & 0.11 & no \\
Llama & E-F1 & $-$0.38 & [$-$0.80, $+$0.05] & 0.08 & no \\
\bottomrule
\end{tabular}
\end{table}

The R-F1 cost is model-dependent: Qwen shows a small but detectable $-$0.51\pp{} decline ($p=0.04$), while Llama shows no significant change ($-$0.30\pp{}, $p=0.11$, CI crosses zero).
Entity F1 changes are not significant for either model.
These results support the ``negligible cost'' characterization: \sys{} composes with fine-tuning at a practically insignificant F1 trade-off while guaranteeing DR below 0.5\%.

\subsection{Cross-Model Generalization}

\sys{} generalizes across model families without code modification.
Table~\ref{tab:cross-model-summary} summarizes DR reduction across all configurations.

\begin{table}[t]
\centering
\small
\caption{Cross-model DR reduction summary (\%).}
\label{tab:cross-model-summary}
\begin{tabular}{llcc}
\toprule
\textbf{Dataset} & \textbf{Regime} & \textbf{Qwen DR red.} & \textbf{Llama DR red.} \\
\midrule
SciERC & Zero-shot & 94.0\% & 99.4\% \\
DocRED & Zero-shot & 44.7\% & 90.2\% \\
DocRED & FT residual & 57.9\% & 69.4\% \\
\bottomrule
\end{tabular}
\end{table}

The cross-model gap on DocRED zero-shot (Qwen 44.7\% vs.\ Llama 90.2\%) is explained by their different RV baselines: Qwen RV$=$29.66\% vs.\ Llama RV$=$16.03\%.
The grammar layer forces OOV relation names into valid vocabulary, revealing latent DR violations that were previously masked by invalid relation names.
This demonstrates that RV and DR are orthogonal failure modes---vocabulary compliance does not imply semantic compliance---and the two constraint layers address them independently.

%==============================================================================
\section{Analysis and Discussion}
\label{sec:analysis}
%==============================================================================

\subsection{Tightness as a Predictive Factor}

Constraint tightness $\tau$ is the primary predictor of residual DR after constrained decoding.
However, DR reduction ceiling is determined by the interaction of tightness with the model's RV baseline.
Models with high RV baselines (Qwen on DocRED: 29.66\%) reveal more latent DR violations after grammar-forced vocabulary correction, yielding lower apparent DR reduction rates than models with low RV baselines (Llama: 16.03\%) on the same ontology.
This is a deeper insight than tightness alone: the same ontology produces different DR reduction ceilings for different models.

\subsection{Constraint Coverage and Fallback}

We define \emph{constraint coverage} as the fraction of type-token positions where semantic masking is active (i.e., $M_{\text{grammar}} \cap M_{\text{semantic}} \neq \emptyset$).
On SciERC ($\bar{\tau}=0.333$), fallback occurs ${\sim}10.7$ times per document; on DocRED ($\bar{\tau}=0.842$), ${\sim}57$ times per document.
This 5.3$\times$ difference directly explains the DR residual gap: positions that fall back to grammar-only retain the grammar layer's baseline DR rate.
FT models produce more well-formed outputs, reducing fallback frequency relative to zero-shot, but the ontology-driven tightness effect persists.

\subsection{Cross-Field Consistency Violations}
\label{sec:cf}

\sys{} trades DR violations for elevated CF violations: the same entity may receive different type labels across different relation contexts.
CF rates are 9.90\% on SciERC (Qwen zero-shot constrained) and 1.50--2.00\% on DocRED (FT+constrained).

\paragraph{Root Cause.}
Ablation confirms CF is introduced entirely by the semantic layer: grammar-only CF$=$0.08\% (matching unconstrained), while full S1 CF$=$9.90\%.
The mechanism is entity-vs-relation type drift (EVR): per-relation independent type enforcement assigns the same entity different types optimized for each relation's constraints.
Over 90\% of CF instances are EVR.

\paragraph{Type System Ambiguity.}
Top conflict patterns---\texttt{Generic}$\leftrightarrow$\texttt{OtherScientificTerm} (60 instances), \texttt{Material}$\leftrightarrow$\texttt{OtherScientificTerm} (28), \texttt{Material}$\leftrightarrow$\texttt{Method} (20)---correspond to genuinely ambiguous type boundaries in SciERC's type system.
The same term can be legitimately classified as different types depending on context.

\paragraph{FT Mitigation.}
Fine-tuning substantially reduces CF: Qwen DocRED CF drops from 2.57\% (zero-shot+\sys{}) to 1.50\% (FT+\sys{}); Llama from 12.02\% to 2.00\%.
FT internalizes entity type conventions, reducing the type drift that semantic enforcement can cause.

\paragraph{Severity Distinction.}
DR violations produce structurally invalid triples unusable by downstream systems.
CF violations produce individually type-valid triples with cross-triple inconsistency---a softer failure mode.
Future work will explore global entity-type consensus passes to resolve CF.

\subsection{S1 vs.\ S2 Architectural Ablation}

We evaluate a two-pass strategy (S2): Pass~1 extracts entities with grammar-only constraints, Pass~2 generates relations with input-dependent grammar compiled from Pass~1 entities.
On SciERC (full evaluation, $n=551$): S2 DR$=$0.97\% (vs.\ S1 0.93\%), CF$=$8.03\% (vs.\ S1 9.90\%), R-F1$=$8.03\% (vs.\ S1 8.36\%).
S2 shows no significant advantage over S1 on any metric and slightly worse F1.
We position S2 as an architectural variant demonstrating framework flexibility rather than a recommended strategy.

%==============================================================================
\section{Related Work}
\label{sec:related}
%==============================================================================

\paragraph{Constrained Decoding for Structured Output.}
Grammar-constrained decoding has been applied to ensure valid JSON, code, and SQL output from LLMs.
XGrammar and related systems compile structural constraints into efficient token masks.
However, these systems enforce purely syntactic constraints and cannot express semantic validity rules such as domain/range type compatibility.
\sys{} extends this paradigm with a semantic layer that enforces cross-field constraints during generation.

\paragraph{LLM-Based Information Extraction.}
Recent work applies LLMs to IE tasks in zero-shot and few-shot settings, achieving competitive performance without task-specific training.
However, the constraint compliance of LLM outputs has not been systematically studied.
Our violation landscape study fills this gap, establishing quantitative baselines for four violation types across multiple models and datasets.

\paragraph{Type-Constrained Knowledge Graph Completion.}
Type constraints have been extensively studied in knowledge graph completion (KGC), where entity type information improves link prediction accuracy.
Unlike KGC, which operates over a fixed entity inventory, IE requires type enforcement during open-ended text generation---a fundamentally different computational setting that motivates our autoregressive constraint approach.

\paragraph{Post-Hoc Validation.}
An alternative to decoding-time enforcement is post-hoc filtering or re-ranking.
Our experiments (\S\ref{sec:posthoc}) show that post-hoc filtering achieves high precision but at severe coverage cost (48.1\% triples rejected on DocRED) and cannot prevent entity-type violations.
\sys{}'s ``prevent rather than detect'' approach preserves information coverage while achieving compliance.

%==============================================================================
\section{Conclusion}
\label{sec:conclusion}
%==============================================================================

We presented a systematic study of ontology constraint violations in LLM-based information extraction, finding that 15--26\% of generated triples violate domain/range type constraints.
Our framework, \sys{}, automatically compiles ontology schemas into composable decoding-time constraints, achieving 90--99\% violation reduction at zero or positive F1 impact in zero-shot settings and negligible F1 cost ($\leq$0.5\pp{}) when composed with fine-tuning.
The two-layer architecture---grammar for structural validity, semantic for type compliance---demonstrates an enabling dependency where grammar constraints are a prerequisite for effective semantic enforcement.

Key findings include: (1)~violation rate is driven by constraint tightness interacting with model vocabulary compliance, not ontology size; (2)~prompt-level intervention is fundamentally unreliable, with some relation types showing \emph{increased} violations; (3)~\sys{} generalizes across model families (Qwen3-8B, Llama-3.1-8B) without code modification; (4)~cross-field consistency violations are an inherent limitation of per-relation independent enforcement, substantially mitigated by fine-tuning.

Future work will address global entity-type consensus to resolve CF violations and extend to additional IE paradigms beyond closed-schema extraction.

\bibliographystyle{plainnat}

\end{document}
